\documentclass[11pt]{article}
\usepackage[T1]{fontenc}
\usepackage{lmodern,amsmath,amssymb,amsthm,mathtools}
\usepackage[margin=1.05in]{geometry}
\usepackage[hidelinks]{hyperref}
\newtheorem{theorem}{Theorem}[section]
\newtheorem{lemma}[theorem]{Lemma}
\newtheorem{proposition}[theorem]{Proposition}
\newtheorem{corollary}[theorem]{Corollary}
\newcommand{\R}{\mathbb R}\newcommand{\E}{\mathbb E}
\newcommand{\Hh}{\mathcal H}\newcommand{\SU}{\mathrm{SU}(2)}
\newcommand{\code}[1]{{\small\nolinkurl{#1}}}
\title{Finite gauge transfer comparisons:\\
seven-link $\mathrm{SU}(2)$ and a homogeneous $\mathbf Z_2$ torus}
\author{Project Stardust}
\date{Reviewed finite-model draft, October 30, 2026}
\begin{document}
\maketitle
\begin{abstract}
We study an explicitly normalized transfer on two adjacent plaquettes
with seven $\mathrm{SU}(2)$ links. Independent central Wilson increments
update each link once, and a spatial half weight is applied on each
side of the kinetic operator. Positive kernel powers construct a
uniformly positive bounded ground function at every fixed coupling,
giving a genuine reversible stationary Markov kernel. For a sufficiently
small positive spatial tilt, depending on the temporal coupling, a
three-channel gauge-invariant trial has normalized stationary lag
strictly larger than the bare one-plaquette character factor. The
proof includes exact Haar projections, a positive supersolution,
explicit cubic errors and a positive optimized second-order margin.
Every fixed number of Markov steps fails to give a parameter-uniform
subunit covariance factor, even after arbitrary random gauge actions.
At each fixed coupling, a separate positive minorization gives geometric
mixing. We also exhibit the loss of the trial's centering symmetry when
both spatial plaquettes are weighted. These comparisons concern finite models;
no homogeneous four-dimensional transfer or physical
mass-gap conclusion is inferred. Classical positive-kernel and character
methods are credited; priority for the exact comparison is unresolved. A separate microscopic homogeneous $\mathbf Z_2$ torus has sharp bare gap $3/4$, but stationary weighted gap at most $39/64$, with a constructed positive algebraic ground function. Separate
conditional scaling and positive Laplace-measure arguments expose the
normalization and spectral hypotheses needed beyond finite transfers.
\end{abstract}

\section{The model and the quantified conclusions}
Identify $G=\SU$ with unit quaternions, and write
$\tau(g)=\frac12\operatorname{Re}\operatorname{tr}(g)$ for its scalar
quaternion coordinate. Let $dg$ be normalized Haar measure. For
$\beta\in\R$ define the actual probability law
\begin{equation}\label{eq:increment}
 Z_\beta=\int_G e^{\beta\tau(v)}\,dv,\qquad
 d\mu_\beta(v)=Z_\beta^{-1}e^{\beta\tau(v)}\,dv,\qquad
 c=c_\beta=\int_G\tau(v)\,d\mu_\beta(v).
\end{equation}
The second scalar channel is
\begin{equation}\label{eq:r}
 m_\beta=\int_G\tau(v)^2\,d\mu_\beta(v),\qquad
 r=r_\beta=\frac{4m_\beta-1}{3}.
\end{equation}
All these quantities are finite, and $Z_\beta>0$, since $|\tau|\leq1$.

Our state space is $\mathcal S=G^7$, with coordinates
$s=(U,A,B,C,D,E,F)$ and product Haar probability $dH(s)$. The six
vertices are numbered $0,\ldots,5$. Orient the shared link $U$ from
$3$ to $0$, the path $(A,B,C)$ from $0$ through $1,2$ to $3$, and
the path $(D,E,F)$ from $3$ through $4,5$ to $0$. Put
\begin{equation}\label{eq:words}
 \begin{gathered}
 L=ABC,\quad R=DEF,\quad x(s)=\tau(LU),\quad y(s)=\tau(U^{-1}R),\\
 z(s)=\frac{\tau(LR)}4,\qquad w(s)=x(s)y(s)-z(s).
 \end{gathered}
\end{equation}
Thus $x,y$ are the two adjacent plaquette traces; $z$ is the shared-link
Haar projection of their product. Every physical link occurs once in
the noise law $\mu_\beta^{\otimes7}$.

On bounded measurable real functions define
\begin{equation}\label{eq:K}
 K_\beta f(s)=\int_{G^7}f(sv)\,d\mu_\beta^{\otimes7}(v),\qquad
 h_t(s)=e^{t x(s)/2},\qquad S_{\beta,t}=M_{h_t}K_\beta M_{h_t}.
\end{equation}
Products $sv$ are coordinatewise. The spatial parameter $t$ and temporal
parameter $\beta$ are independent. Only the first plaquette is weighted
in \eqref{eq:K}. Section~\ref{sec:homogeneous} explains why that
qualification cannot be removed from the slow-mode proof.

\begin{theorem}\label{thm:main}
For each $\beta,t\in\R$, there are $\Lambda>0$ and a bounded measurable
$\psi:\mathcal S\to(0,\infty)$, bounded away from zero, such that
$S_{\beta,t}\psi=\Lambda\psi$. They define a reversible Markov kernel
\begin{equation}\label{eq:P}
 P_{\beta,t}f(s)=\frac{S_{\beta,t}(\psi f)(s)}{\Lambda\psi(s)},
 \qquad d\pi_{\beta,t}=\frac{\psi^2\,dH}{\int\psi^2\,dH}.
\end{equation}
The ground function can be chosen even under the private-link center
flip $D\mapsto-D$ and invariant under the six-vertex gauge group.

If $c_\beta\ne0$, there is $\delta_\beta\in(0,1]$ such that for every
$0<t<\delta_\beta$ a bounded gauge-invariant $f$ satisfies
\begin{equation}\label{eq:mainlag}
 \int f\,d\pi_{\beta,t}=0,\qquad \int f^2\,d\pi_{\beta,t}>0,
 \qquad c_\beta^4<
 \frac{\int fP_{\beta,t}f\,d\pi_{\beta,t}}{\int f^2\,d\pi_{\beta,t}}.
\end{equation}
For the same $f$, every integer $n\geq0$ obeys
\begin{equation}\label{eq:lagiter}
 1-n(1-c_\beta^4)\leq
 \frac{\int fP_{\beta,t}^{\,n}f\,d\pi_{\beta,t}}{\int f^2\,d\pi_{\beta,t}}
 \leq1.
\end{equation}
For each fixed $n$, $q<1$ and $B\in\R$, one can choose
$\beta>B$ and $0<t\leq1$ with a centered gauge-invariant $f$ whose
normalized $n$-step lag exceeds $q$. These conclusions persist if an
arbitrary probability-distributed gauge action is applied before every
Markov step.
\end{theorem}

This is a statement about actual finite-dimensional probability laws
and centered covariances. Its positive and negative conclusions concern
different quantifiers: fixed-parameter mixing versus uniform contraction
over parameters. The proof will retain $\delta_\beta$ throughout.
Positive integral operators have a classical Perron--Jentzsch theory,
including Birkhoff's projective treatment \cite{Birkhoff}. Our elementary
power construction makes the required input explicit. L\"uscher's
physical transfer construction \cite{Luscher} additionally identifies a
gauge-invariant quantum state space and a lattice Hamiltonian; those
identifications are not asserted for this seven-link Markov model.

\section{Haar coordinates and the three exact channels}
If $V=(a,b)\in\R\times\R^3$ is a unit quaternion, then $\tau(V)=a$.
Uniform measure on the unit sphere in $\R^4$ is Haar measure on $G$.
Its coordinate symmetries give
\begin{equation}\label{eq:Haarmoments}
 \int a\,dV=0,\quad \int a^2\,dV=\frac14,\quad
 \int a^4\,dV=\frac18,\qquad
 \int a_i a_j\,dV=\frac{\delta_{ij}}4.
\end{equation}
For the fourth moment, rotational symmetry gives
$\E a_i^4=3\E a_i^2a_j^2$ for $i\ne j$, by rotating
$(a_i+a_j)/\sqrt2$. Expanding $(\sum_{i=0}^3a_i^2)^2=1$
then gives $4\E a_0^4+12\E a_0^2a_1^2=1$ and the stated value.

The tilt \eqref{eq:increment} depends only on $a$, so
\begin{equation}\label{eq:tiltmoments}
 \E_{\mu_\beta}b_i=\E_{\mu_\beta}ab_i=0,\qquad
 \E_{\mu_\beta}b_i b_j=\frac{1-m_\beta}{3}\delta_{ij}.
\end{equation}
Consequently, for each fixed $Q\in G$,
\begin{align}
 \int\tau(QV)\,d\mu_\beta(V)&=c\tau(Q),\label{eq:linear}\\
 \int[\tau(QV)^2-1/4]\,d\mu_\beta(V)
 &=r[\tau(Q)^2-1/4].\label{eq:quadratic}
\end{align}
Indeed $\tau(QV)=a_Qa-b_Q\cdot b$, and
\eqref{eq:tiltmoments} evaluates both expressions. More generally,
right multiplication by the tilted quaternion contracts every linear
form in the four quaternion coordinates by $c$, and every traceless
quadratic form by $r$. For the latter assertion, let $O_V$ be the
orthogonal matrix of right multiplication. Averaging
$O_V^TTO_V$ for a symmetric matrix $T$ gives
$rT+(1-r)\operatorname{tr}(T)I/4$.
This follows by writing $O_V=aI+\sum b_iJ_i$, using the vanishing
mixed moments and
$\sum_{i=1}^3J_i^TTJ_i=\operatorname{tr}(T)I-T$ for symmetric $T$;
the identity is checked on diagonal and symmetric off-diagonal matrix
units using the quaternion multiplication table.

\begin{lemma}\label{lem:channelbounds}
For every finite $\beta$,
$|c|<1$, $|r|<1$, and $c^2\leq m_\beta<1$.
Also $c_\beta\to1$ as $\beta\to+\infty$.
\end{lemma}
\begin{proof}
Cauchy--Schwarz gives $c^2\leq m_\beta$. The tilted measure has a
strictly positive density relative to Haar measure, and $a^2<1$ on a
nonempty open set, so $m_\beta<1$. Hence $|c|<1$ and
$-1/3\leq r<1$, implying $|r|<1$.
For $\varepsilon>0$, the upper cap $\{a>1-\varepsilon/2\}$ has
Haar mass $b_\varepsilon>0$, and comparison of numerator and partition
function gives, for $\beta\geq0$,
\[
 \mu_\beta\{a\leq1-\varepsilon\}
 \leq b_\varepsilon^{-1}e^{-\beta\varepsilon/2}.
\]
Thus $0\leq1-c_\beta\leq\varepsilon+
2b_\varepsilon^{-1}e^{-\beta\varepsilon/2}$.
First let $\beta\to\infty$ and then $\varepsilon\downarrow0$.
\end{proof}

Under product Haar measure, the three variables $L,R,U$ are independent
Haar variables. Each path is a product of independent Haar links, and
the two paths and shared link use disjoint coordinates. Quaternion
coordinate integration gives
\begin{equation}\label{eq:pairtrace}
 \int_G\tau(PV)\tau(QV)\,dV=\frac{\tau(PQ^{-1})}{4}.
\end{equation}
Changing variables $V=LU$ in this formula proves
$\int_G x(s)y(s)\,dU=z(s)$ for fixed $L,R$.

\begin{proposition}\label{prop:channels}
The functions $y,z,w$ are orthogonal in $L^2(H)$. Their squared
norms are $1/4$, $1/64$ and $3/64$, respectively. Their exact kinetic
evolutions are
\begin{equation}\label{eq:threeeigen}
 K_\beta y=c^4y,\qquad K_\beta z=c^6z,\qquad
 K_\beta w=c^6r\,w.
\end{equation}
\end{proposition}
\begin{proof}
For fixed $L,U$, write $x=\tau(LU)$. Integration in the independent
right-path Haar variable $R$ gives
\[
 \E_R y^2=\tfrac14,\qquad \E_R z^2=\tfrac1{64},\qquad
 \E_R yz=\tfrac{x}{16}.
\]
Since $w=xy-z$, this also gives
$\E_R yw=3x/16$, $\E_R zw=x^2/16-1/64$ and
$\E_R w^2=x^2/8+1/64$. Averaging over $L,U$ and using
$\E x=0$, $\E x^2=1/4$ proves all norms and cross products.

Each of the four independent link increments in the right plaquette
contributes the factor $c$ in \eqref{eq:linear}; inverses have the same
law, so this proves the first evolution. The six outer increments in
$\tau(LR)$ similarly give $c^6$. For the last evolution, first
integrate the six outer increments. The product $xy$ acquires $c^6$
and so does $z$. As a function of the remaining shared increment,
$xy-z$ is a quadratic form in its quaternion coordinates with zero
Haar mean. It is therefore a traceless quadratic form, contracted by
$r$ according to the calculation following \eqref{eq:quadratic}.
This shared link is integrated once, giving $c^6r$ rather than a
factor obtained by pretending the two plaquettes have independent
shared links.
\end{proof}

The law $\mu_\beta$ is invariant under inversion and conjugation.
Thus $K_\beta$ preserves Haar measure and is selfadjoint on $L^2(H)$,
by the change of variables $(s,v)\mapsto(sv,v^{-1})$ in the bounded
pairing. Jensen's inequality gives its $L^2$ contraction; bounded
functions are dense, so these statements extend to the whole space.
The operator $S_{\beta,t}$ is bounded and selfadjoint. Positive
integral kernels below mean pointwise positive densities, not an
unproved assertion that every $L^2$ eigenvalue is nonnegative.

\section{A ground function constructed from positive kernel powers}
We give the construction because a conditional ground-state formula
cannot alone establish the stationary assertions in Theorem~\ref{thm:main}.

\begin{lemma}\label{lem:positivepowers}
Let $(X,\nu)$ be a nonempty probability space and let $k:X\times X\to\R$
be jointly measurable with $0<m\leq k(x,y)\leq M<\infty$ everywhere.
Then $Tf(x)=\int k(x,y)f(y)\,d\nu(y)$ has a measurable eigenfunction
$\psi$ and eigenvalue $\Lambda$ satisfying
\[
 m\leq\Lambda\leq M,\qquad m/M\leq\psi(x)\leq M/m,
 \qquad T\psi=\Lambda\psi.
\]
The construction requires no supplied eigenfunction or spectral radius.
\end{lemma}
\begin{proof}
Set $p_0=1$, $p_{n+1}=Tp_n$ and $r_n=p_{n+1}/p_n$.
Then $m^n\leq p_n\leq M^n$, and positivity gives
$m\leq r_n\leq M$ by induction: $mp_n\leq p_{n+1}\leq Mp_n$
holds at $n=0$ and is preserved by $T$.
Let $\nu_n=p_n\nu/\int p_n\,d\nu$. The exact ratio relation is
\[
 r_{n+1}(x)=\frac{\int k(x,y)r_n(y)\,d\nu_n(y)}
                         {\int k(x,y)\,d\nu_n(y)}.
\]
The averaging density in this expression is at least
$\varepsilon=m/M$. Subtracting $\varepsilon\nu_n$ leaves a positive
measure of mass $1-\varepsilon$. If $r_n\in[a_n,b_n]$, it follows that
\[
 r_{n+1}\in[a_{n+1},b_{n+1}],\qquad
 (a_{n+1},b_{n+1})=
 \varepsilon\!\int r_n\,d\nu_n+(1-\varepsilon)(a_n,b_n).
\]
Starting with $[a_0,b_0]=[m,M]$, these intervals are nested and
$b_n-a_n=(M-m)(1-\varepsilon)^n$. They shrink to one common
$\Lambda\in[m,M]$, so $r_n\to\Lambda$ uniformly.

Fix $x_0\in X$ and set $q_n(x)=p_n(x)/p_n(x_0)$.
The comparison $m\int p_{n-1}\leq p_n(x)\leq M\int p_{n-1}$ gives
$m/M\leq q_n\leq M/m$ for $n\geq1$, and also for $n=0$.
Since $q_{n+1}(x)=q_n(x)r_n(x)/r_n(x_0)$, we have
\[
 |q_{n+1}(x)-q_n(x)|
 \leq\frac{M(M-m)}{m^2}(1-\varepsilon)^n.
\]
The geometric series proves uniform convergence to a measurable
$\psi$ with the stated bounds. The identity
$Tq_n=r_n(x_0)q_{n+1}$ passes to the limit by dominated convergence,
giving $T\psi=\Lambda\psi$ pointwise. The case $m=M$ is included:
the kernel is constant and all subsequent ratios and normalized powers
are constant.
\end{proof}

For \eqref{eq:K}, the actual symmetric density relative to $H$ is
\begin{equation}\label{eq:density}
 k_{\beta,t}(s,r)=h_t(s)h_t(r)
 Z_\beta^{-7}\exp\!\left\{\beta\sum_{e=1}^7\tau(s_e^{-1}r_e)\right\}.
\end{equation}
It satisfies the everywhere bounds
\begin{equation}\label{eq:kernelbounds}
 m_{\beta,t}=\frac{e^{-|t|-7|\beta|}}{Z_\beta^7}
 \leq k_{\beta,t}\leq
 \frac{e^{|t|+7|\beta|}}{Z_\beta^7}=M_{\beta,t}.
\end{equation}
Lemma~\ref{lem:positivepowers} therefore proves actual ground-function
existence for every real $\beta,t$. Formula \eqref{eq:P} is a genuine
probability kernel, since its density integrates to one by the
eigenfunction equation. Detailed balance follows because
\[
 \pi(ds)P(s,dr)=
 \frac{\psi(s)k_{\beta,t}(s,r)\psi(r)}
      {\Lambda\int\psi^2\,dH}\,dH(s)dH(r)
\]
is symmetric. In particular $\pi$ is stationary, and $P$ is a
selfadjoint contraction on the full $L^2(\pi)$ space.

\section{Symmetries, centering and the actual observable}
For $g=(g_0,\ldots,g_5)\in G^6$, the gauge action on an oriented edge
$i\to j$ is $s_e\mapsto g_i s_e g_j^{-1}$. The words $LU$, $U^{-1}R$
and $LR$ are closed words based at vertex $0$, so their half traces,
and hence $x,y,z,w$, are gauge invariant. Product Haar measure is
invariant under this action. Simultaneous transformation of $s,r$ in
\eqref{eq:density} conjugates each $s_e^{-1}r_e$ by the gauge element
at its terminal vertex; $\tau$ is conjugation invariant. Thus
\begin{equation}\label{eq:gaugecovariance}
 k_{\beta,t}(gs,gr)=k_{\beta,t}(s,r),\qquad
 S(f\circ g)=(Sf)\circ g.
\end{equation}

Let $J$ be the center flip $D\mapsto-D$ on the first private link of
the right path. Then $J^2=I$, $H$ is $J$-invariant, and
\[
 x\circ J=x,\qquad (y,z,w)\circ J=-(y,z,w).
\]
The operator $K_\beta$ commutes with this left multiplication, since
$J(s)v=J(sv)$; its two spatial weights are also unchanged.
Replacing a positive ground function by $(\psi+\psi\circ J)/2$
therefore preserves its eigenvalue, lower and upper bounds, and makes
it $J$-even. This replacement is made from now on.

\begin{lemma}\label{lem:gaugeground}
Every uniformly positive bounded measurable eigenfunction of $S$
is gauge invariant. In particular the chosen even $\psi$, the stationary
measure $\pi$ and its kernel $P$ have the stated gauge covariance.
\end{lemma}
\begin{proof}
Fix a gauge transformation $g$. The function $\chi=\psi\circ g$
has the same eigenvalue by \eqref{eq:gaugecovariance}.
The ratio $F=\chi/\psi$ is bounded and satisfies $PF=F$.
The ground transform can be written as normalized averaging of the
same kernel $k$ against the probability
\begin{equation}\label{eq:bridgereference}
 d\nu_\psi=\frac{\psi\,dH}{\int\psi\,dH}:
 \quad Pf(s)=\frac{\int k(s,r)f(r)\,d\nu_\psi(r)}
                       {\int k(s,r)\,d\nu_\psi(r)}.
\end{equation}
The averaging density is at least $m_{\beta,t}/M_{\beta,t}=\varepsilon>0$.
If $F$ lies in $[a,b]$, subtracting the common
$\varepsilon\nu_\psi$ part shows that the range of $PF$ has width
at most $(1-\varepsilon)(b-a)$. Since $PF=F$, the oscillation of $F$
is zero, and $\chi=C\psi$ pointwise. Gauge invariance of $H$ gives
$\int\chi\,dH=\int\psi\,dH>0$, forcing $C=1$.
The density formula for $\pi$ and \eqref{eq:gaugecovariance} now prove
invariance of the measure and covariance of the kernel. Notice that
$\nu_\psi$ in \eqref{eq:bridgereference} is generally different from
the stationary measure with density proportional to $\psi^2$.
\end{proof}

For real $A,B$ define the concrete odd trial
\begin{equation}\label{eq:oddtrial}
 p_t=y+tAz+tBw,\qquad F_t=h_t^{-1}p_t,
 \qquad f_t=F_t/\psi.
\end{equation}
The trial and observable are bounded and measurable. The lower bound
on $\psi$ is needed for boundedness of $f_t$. They are gauge invariant,
$F_t$ is $J$-odd, and $\psi$ is $J$-even. Therefore
\begin{equation}\label{eq:center}
 \int f_t\,d\pi=\frac{\int F_t\psi\,dH}{\int\psi^2\,dH}=0.
\end{equation}
The exact norm and lag relation is
\begin{equation}\label{eq:rayleigh}
 \frac{\int f_tPf_t\,d\pi}{\int f_t^2\,d\pi}
 =\frac{\langle F_t,SF_t\rangle_H}{\Lambda\|F_t\|_H^2}.
\end{equation}
Both numerator and denominator follow by cancelling the positive
$\psi$ factors in \eqref{eq:P}. This equality uses the actual ground
normalization; a pointwise row normalization of $S$ would define a
different operator and would not give this formula.

\section{A positive supersolution with a cubic error}
We first bound $\Lambda$ above on a nonempty interval of spatial
parameters. This will supply the normalization in the lower trial test.
Set
\begin{equation}\label{eq:vacuumconstants}
 a=c^4,\quad b=r^4,\qquad
 u=\frac{a}{1-a},\quad
 v=\frac{b(1+a)}{2(1-a)(1-b)},\quad
 \mu=\frac{1+a}{8(1-a)}.
\end{equation}
The denominators are nonzero by Lemma~\ref{lem:channelbounds}.
For $|q|\leq1$ let
\begin{align}
 P_t(q)&=1+utq+vt^2(q^2-1/4),\label{eq:polynomial}\\
 C_0(A,B)&=(2+e)(1+|A|+|B|),\nonumber\\
 C&=C_0(u,v)+|\mu|(|u|+|v|),\qquad
 Q_t=1+\mu t^2+2C|t|^3.\label{eq:certificate}
\end{align}

\begin{lemma}\label{lem:taylor}
For $|t|\leq1$ and $|q|\leq1$,
\begin{equation}\label{eq:taylor}
 \left|e^{tq}(1+Atq+Bt^2(q^2-1/4))
 -\left[1+(1+A)tq+
 t^2\{(B+A+1/2)q^2-B/4\}\right]\right|
 \leq C_0(A,B)|t|^3.
\end{equation}
If a probability-distributed $Y$ takes values in $[-1,1]$ and has
$\E Y=aq$, $\E Y^2=1/4+b(q^2-1/4)$, then
\begin{equation}\label{eq:supersolutionerror}
 |\E[e^{tY}P_t(Y)]-(1+\mu t^2)P_t(q)|\leq C|t|^3.
\end{equation}
If $|u||t|+|v||t|^2\leq1/2$, then $P_t(q)\geq1/2$ and
$\E[e^{tY}P_t(Y)]\leq Q_tP_t(q)$.
\end{lemma}
\begin{proof}
Taylor's series gives
$|e^{tq}-1-tq-(tq)^2/2|\leq e|t|^3$ on the stated rectangle.
Multiplication by $1+Atq+Bt^2(q^2-1/4)$ has absolute value at most
$1+|A|+|B|$, since $|q^2-1/4|\leq1$. The remaining terms of total
degree three and four are
\[
 t^3[Bq(q^2-1/4)+Aq^3/2]+t^4Bq^2(q^2-1/4)/2,
\]
bounded by $|t|^3(3|B|/2+|A|/2)$. This proves
\eqref{eq:taylor} with the displayed larger constant.

Substitute $A=u,B=v$ and the two moments of $Y$. The definitions
$a(1+u)=u$ and $b(v+u+1/2)=v$ give exactly
\[
 \E\left[1+(1+u)tY+t^2\{(v+u+1/2)Y^2-v/4\}\right]
 =P_t(q)+\mu t^2.
\]
Also $|P_t(q)-1|\leq|t|(|u|+|v|)$ for $|t|\leq1$.
The difference between $P_t+\mu t^2$ and $(1+\mu t^2)P_t$
is therefore bounded by $|\mu|(|u|+|v|)|t|^3$, proving
\eqref{eq:supersolutionerror}. The smallness condition gives
$P_t\geq1/2$, which absorbs the additive $C|t|^3$ into
$2C|t|^3P_t$ and proves the last assertion.
\end{proof}

\begin{proposition}\label{prop:lambda}
If $|t|\leq1$ and $|u||t|+|v||t|^2\leq1/2$, then
$\Lambda\leq Q_t$.
\end{proposition}
\begin{proof}
For $q=x(s)$ the random variable $Y=x(sV)$ under the seven-link
increment law has the moments in Lemma~\ref{lem:taylor}: the four
increments in the first plaquette give $a=c^4$ and $b=r^4$ by
\eqref{eq:linear}--\eqref{eq:quadratic}. Inverses and conjugations
preserve the central increment law, so these moment computations are
valid for every fixed starting configuration. Thus
$U_t(s)=h_t(s)P_t(x(s))$ satisfies $SU_t\leq Q_tU_t$.
It is bounded and at least $e^{-|t|/2}/2$.

For any bounded measurable $R$, symmetry and positivity of $k$ give
\begin{align*}
 \langle U_tR,SU_tR\rangle_H
 &\leq\frac12\iint k(s,r)U_t(s)U_t(r)[R(s)^2+R(r)^2],dH(s)dH(r)\\
 &=\int U_t(s)(SU_t)(s)R(s)^2\,dH(s)
 \leq Q_t\|U_tR\|_H^2.
\end{align*}
The first inequality is $2R(s)R(r)\leq R(s)^2+R(r)^2$.
Take $R=\psi/U_t$, which is bounded by the positive lower bound
on $U_t$. The eigenfunction equation and $\|\psi\|_H>0$ give
$\Lambda\leq Q_t$. Neither a formal eigenbranch nor a spectral
identification of the approximate polynomial is used.
\end{proof}

\section{The exact odd numerator and denominator}
The inverse half weight in \eqref{eq:oddtrial} cancels the two spatial
weights in the kinetic numerator. Proposition~\ref{prop:channels} gives
\begin{equation}\label{eq:numerator}
 N_t:=\langle F_t,SF_t\rangle_H
 =\frac{c^4}4+\frac{t^2}{64}(c^6A^2+3c^6rB^2).
\end{equation}
There is no remainder in this identity. All exponential error lies in
the denominator $D_t=\|F_t\|_H^2=\int e^{-tx}p_t^2\,dH$.

\begin{lemma}\label{lem:denominator}
For all real $t,A,B$, one has $D_t\geq e^{-|t|}/4>0$ and
\begin{equation}\label{eq:radial}
 D_t=\int_G e^{-t\tau(V)}\mathcal R_t(\tau(V))\,dV,
 \quad \mathcal R_t(q)=\frac14+\frac{t(A+3B)q}{8}
 +\frac{t^2}{64}[(A-B)^2+8B(A+B)q^2].
\end{equation}
Put
\begin{align}
 d&=\frac{A^2+3B^2-2A-6B+2}{64},\nonumber\\
 R_0&=\frac{(A-B)^2}{64}+\frac{B(A+B)}{32},\nonumber\\
 E&=\frac14C_0(-(A+3B)/2,B(A+B)/2)+2|R_0|.\label{eq:Derror}
\end{align}
Then, for $|t|\leq1$,
\begin{equation}\label{eq:Dexpansion}
 |D_t-(1/4+dt^2)|\leq E|t|^3.
\end{equation}
\end{lemma}
\begin{proof}
For fixed $L,U$, $x=\tau(LU)$ and
$p_t=(1+tBx)y+t(A-B)z$. The three conditional moments in the proof
of Proposition~\ref{prop:channels} therefore give
$\E_Rp_t^2=\mathcal R_t(x)$. Haar invariance in $L$ makes $LU$
a Haar variable, proving \eqref{eq:radial}.
The exact global Gram identity also gives
$\int p_t^2\,dH=1/4+t^2(A^2+3B^2)/64\geq1/4$;
since $e^{-tx}\geq e^{-|t|}$, the positive lower bound follows.

Set $A_0=-(A+3B)/2$, $B_0=B(A+B)/2$. Direct expansion shows
\[
 \mathcal R_t(q)=\tfrac14[1+A_0(-t)q+B_0(-t)^2(q^2-1/4)]
                  +t^2R_0.
\]
Apply \eqref{eq:taylor} with parameter $-t$ to the bracket.
In the last term use $|e^{-tq}-1|\leq2|t|$ for $|t|,|q|\leq1$.
The total error is at most $E|t|^3$. The integrated linear part
vanishes by $\int\tau=0$; the quadratic part is $dt^2$ by
$\int\tau^2=1/4$. This proves \eqref{eq:Dexpansion} with the stated
constant and both signs of $t$.
\end{proof}

\section{An optimized positive margin and a genuine spatial interval}
Assume $c\ne0$ and put $\alpha=c^2$. Fix the actual trial parameters
\begin{equation}\label{eq:AB}
 A=\frac1{1-\alpha},\qquad B=\frac1{1-\alpha r}.
\end{equation}
Their denominators are strictly positive by Lemma~\ref{lem:channelbounds}.
Let $d,E$ be \eqref{eq:Derror} for this pair, and define
\begin{equation}\label{eq:margin}
 \Delta=\frac{c^6A^2+3c^6rB^2}{64}-c^4d-\frac{c^4\mu}{4}.
\end{equation}

\begin{lemma}\label{lem:margin}
The margin is strictly positive. More explicitly,
\begin{equation}\label{eq:marginfactor}
 \Delta=\frac{\alpha^3[1+3r-\alpha(3+r)]}
 {64(1-\alpha^2)(1-\alpha r)}>0.
\end{equation}
\end{lemma}
\begin{proof}
Insert \eqref{eq:AB} and $\mu=(1+c^4)/(8(1-c^4))$ into
\eqref{eq:margin}. Completing the two scalar squares in $A,B$ gives
\[
 \Delta=\frac{\alpha^2}{64}
 \left[\frac1{1-\alpha}+\frac3{1-\alpha r}
                    -\frac4{1-\alpha^2}\right].
\]
Putting the bracket over its common denominator yields
\eqref{eq:marginfactor}. Now $0<\alpha\leq m_\beta<1$ and
$1+3r=4m_\beta\geq4\alpha$, while $r<1$. Thus
$1+3r-\alpha(3+r)\geq\alpha(1-r)>0$.
Both denominator factors are positive, proving the assertion.
\end{proof}

This margin compares the exact odd trial with an upper bound on the
actual ground normalization. A calculation of the kinetic numerator
alone would not establish \eqref{eq:mainlag}.
Define the nonnegative comparison constant
\begin{equation}\label{eq:comparisonconstant}
 J=E+\mu(|d|+E)+2C(1/4+|d|+E),
\end{equation}
and choose the explicit positive window
\begin{equation}\label{eq:delta}
 \delta_\beta=\min\left\{1,\frac{\Delta}{|c^4J|+1},
                     \frac1{2(|u|+|v|+1)}\right\}.
\end{equation}
For $0<t<\delta_\beta$, both supersolution smallness and
$c^4Jt<\Delta$ hold. From \eqref{eq:Dexpansion}, for $0\leq t\leq1$,
\[
 D_t\leq\tfrac14+|d|+E,\qquad
 D_t-\tfrac14\leq t(|d|+E).
\]
Expanding $Q_tD_t=D_t+\mu t^2D_t+2Ct^3D_t$ and using these two
bounds gives
\begin{equation}\label{eq:QD}
 Q_tD_t\leq\tfrac14+(d+\mu/4)t^2+Jt^3.
\end{equation}
Together with \eqref{eq:numerator} and \eqref{eq:margin}, this yields
\[
 N_t-c^4Q_tD_t\geq t^2(\Delta-c^4Jt)>0.
\]
Proposition~\ref{prop:lambda} gives $\Lambda\leq Q_t$, while
$D_t,\Lambda>0$. Equation \eqref{eq:rayleigh} therefore proves
\eqref{eq:mainlag}. Centering was already proved in \eqref{eq:center},
and positive variance follows from $D_t>0$. This proves the one-step
part of Theorem~\ref{thm:main} on an actual nonempty interval, with no
unquantified perturbative error. The window may shrink as $\beta$ grows.

\section{Finite-step persistence and gauge randomization}
We need an iteration argument for the actual stationary operator,
including a possible negative part of its spectrum.

\begin{lemma}\label{lem:iteration}
Let $T$ be a selfadjoint contraction on a real Hilbert space, and set
$\mathcal D(F)=\|F\|^2-\langle TF,F\rangle$. Then, for all $n\geq0$,
\begin{equation}\label{eq:Hilbertiteration}
 \|F\|^2-n\mathcal D(F)\leq\langle T^nF,F\rangle\leq\|F\|^2.
\end{equation}
\end{lemma}
\begin{proof}
Cauchy--Schwarz and contraction give $\mathcal D\geq0$ and
$\langle TZ,Z\rangle\geq-\|Z\|^2$.
For $Z=F-TF$, expansion using selfadjointness gives
\[
 \mathcal D(F)-\mathcal D(TF)
 =\langle (I+T)Z,Z\rangle\geq0.
\]
Also $\|F\|^2-\|TF\|^2\leq2\mathcal D(F)$, because the right
side minus the left side is $\|F-TF\|^2$.
Iterating yields $\|T^kF\|^2\geq\|F\|^2-2k\mathcal D(F)$.
For $n=2k$, selfadjointness gives
$\langle T^{2k}F,F\rangle=\|T^kF\|^2$.
For $n=2k+1$ it gives
\[
 \langle T^{2k+1}F,F\rangle
 =\|T^kF\|^2-\mathcal D(T^kF)
 \geq\|F\|^2-(2k+1)\mathcal D(F).
\]
These are the lower bounds. The upper bound follows from
$\|T^nF\|\leq\|F\|$ and Cauchy--Schwarz. No nonnegative-spectrum
premise has entered the proof.
\end{proof}

Apply the lemma with $T=P_{\beta,t}$ on $L^2(\pi_{\beta,t})$ and
the observable in \eqref{eq:oddtrial}. Equation \eqref{eq:mainlag}
gives $\mathcal D(f)<(1-c_\beta^4)\|f\|^2$, proving
\eqref{eq:lagiter}. Fix an integer $n\geq0$, $q<1$ and a cutoff $B$.
For $n\geq1$, Lemma~\ref{lem:channelbounds} allows $\beta>B$
so large that $1-n(1-c_\beta^4)>q$ and $c_\beta\ne0$.
Choose $t=\delta_\beta/2$. The $n$-step ratio exceeds $q$.
For $n=0$ it is one, and the same construction supplies a nonzero
centered invariant observable. This proves nonuniformity without
passing to any continuum or volume limit.

Let $\nu$ be any probability measure on $G^6$ and define the actual
randomized Markov operator
\begin{equation}\label{eq:gaugeaverage}
 \overline P_\nu f(s)=\int_{G^6}(Pf)(gs)\,d\nu(g).
\end{equation}
This randomizes the starting state by a gauge action before one $P$
step. Gauge invariance of $\pi$ and stationarity for $P$ give
$\pi\overline P_\nu=\pi$. If $f$ is invariant, covariance
\eqref{eq:gaugecovariance} makes $Pf$ invariant too; hence
$\overline P_\nu f=Pf$. Induction gives
\begin{equation}\label{eq:gaugeiterate}
 \overline P_\nu^{\,n}f=P^nf\qquad(n\geq0)
\end{equation}
for every bounded invariant $f$. All the covariance ratios are
unchanged, including for product Haar randomization at all six
vertices. The randomized operator need not itself be selfadjoint:
its action on the invariant observables has already been identified
with the selfadjoint $P$ iterates. This completes
Theorem~\ref{thm:main}.

A single open link would give a misleading alternative test. If
$F:G\to\R$ satisfies $F(gAh^{-1})=F(A)$ for all independent endpoint
gauges, setting $A=1,h=1$ gives $F(g)=F(1)$, so $F$ is constant.
The half trace of an open link is not invariant: its values at $1$
and $-1$ are $1$ and $-1$. The observables above use actual closed
words and do not rely on an open-link mode being physical.

\section{Mixing at each fixed coupling}
The absence of a parameter-uniform fixed-step factor does not prevent
mixing at a fixed pair of parameters.

\begin{proposition}\label{prop:mixing}
For every fixed $\beta,t$ and bounded measurable $f$ with
$\int f\,d\pi_{\beta,t}=0$ and $|f|\leq B$,
\begin{equation}\label{eq:mixing}
 |P_{\beta,t}^{\,n}f(s)|\leq2B(1-\varepsilon_{\beta,t})^n,
 \qquad\varepsilon_{\beta,t}=e^{-2|t|-14|\beta|}>0.
\end{equation}
\end{proposition}
\begin{proof}
Use the probability $\nu_\psi$ in \eqref{eq:bridgereference}.
The normalized row density of $k_{\beta,t}$ is at least
$m_{\beta,t}/M_{\beta,t}=\varepsilon_{\beta,t}$ by
\eqref{eq:kernelbounds}; all seven Wilson partitions cancel.
Thus $P$ is the sum of a common reset
$\varepsilon_{\beta,t}\nu_\psi$ and a positive residual kernel of
mass $1-\varepsilon_{\beta,t}$. It follows that
\[
 \operatorname{osc}(Pf)\leq
 (1-\varepsilon_{\beta,t})\operatorname{osc}(f),
 \quad \operatorname{osc}(f)=\sup f-\inf f.
\]
Iteration bounds the oscillation by $2B(1-\varepsilon_{\beta,t})^n$.
Stationarity gives $\int P^nf\,d\pi=0$, so zero lies between its
infimum and supremum. Its absolute value is bounded by its oscillation,
proving \eqref{eq:mixing}. If $\varepsilon=1$, the kernel is a reset;
the estimate for $n\geq1$ is direct and the $n=0$ bound is immediate.
\end{proof}

This is a bounded-observable mixing result for the actual stationary
finite model. Its reset constant degenerates as $\beta$ grows. The
slow-mode theorem proves failure of a uniform fixed-step factor,
rather than merely pointing to this weak upper estimate.

\section{Why the second spatial plaquette changes the argument}
\label{sec:homogeneous}
To weight both plaquettes equally, replace $h_t$ by
\[
 \widetilde h_t(s)=e^{t(x(s)+y(s))/2},\qquad
 \widetilde S_{\beta,t}=M_{\widetilde h_t}K_\beta M_{\widetilde h_t}.
\]
The old center flip sends $x+y$ to $x-y$, so it no longer preserves
the weight at $t\ne0$. At the flat configuration all links are $1$;
the two values of $\widetilde h_t$ before and after the flip are
$e^t$ and $1$. This is a pointwise failure of a required symmetry,
which persists for every nonzero $t$.

The effect on centering is realized at the genuine reset coupling
$\beta=0$. Here $\mu_0$ is Haar and $K_0f=\int f\,dH$, so
\[
 \widetilde S_{0,t}f=\widetilde h_t\int\widetilde h_tf\,dH,
 \quad\widetilde\psi_t=\widetilde h_t,
 \quad\widetilde\Lambda_t=\int e^{t(x+y)}\,dH,
 \quad d\widetilde\pi_t=
 \frac{e^{t(x+y)}\,dH}{\int e^{t(x+y)}\,dH}.
\]
The actual ground-normalized Markov kernel is the constant reset
$\widetilde P_t f=\int f\,d\widetilde\pi_t$.
Bounded differentiation under the integral gives
\begin{equation}\label{eq:rightmean}
 \left.\frac{d}{dt}\int y\,d\widetilde\pi_t\right|_{t=0}
 =\int y(x+y)\,dH-\int y\,dH\int(x+y)\,dH=\frac14.
\end{equation}
All first moments vanish, $\int xy\,dH=0$ by integrating the
independent left Haar path, and $\int y^2\,dH=1/4$.
The mean at zero is zero; \eqref{eq:rightmean} proves positive mean
for arbitrarily small positive $t$. Under these actual stationary laws,
\[
 \int y\widetilde P_ty\,d\widetilde\pi_t
 =\left(\int y\,d\widetilde\pi_t\right)^2>0,
\]
whereas every correctly centered integrable observable has exactly
zero one-step lag under the reset kernel. An uncentered positive lag
therefore does not obstruct mixing.

This example does not rule out a corrected two-weight trial. It shows
why the one-weight orthogonality and its covariance conclusion cannot
simply be transplanted. In particular, the small independent $t$
interval in Theorem~\ref{thm:main} gives no statement along a
homogeneous relation between spatial and temporal couplings.

\section{A separate homogeneous finite $\mathbf Z_2$ comparison}
The two-weight $\SU$ question is not settled by changing the group.
Nevertheless the frozen library contains a useful independent test of
an unchanged-gap claim under genuinely homogeneous Wilson weights.
We reconstruct it explicitly. This section uses the periodic two-by-two
$\mathbf Z_2$ lattice, not the seven-link state space.

\subsection*{The microscopic law and its actual gauge quotient}
Represent the eight link signs by $u\in\mathbf F_2^8$, with link $j$
equal to $(-1)^{u_j}$. The four vertex-flip incidence masks have supports
\[
 \{0,1,3,4\},\quad\{1,2,3,6\},\quad\{0,4,5,7\},
 \quad\{2,5,6,7\}.
\]
Their sum is zero. The first three are independent, so the effective
gauge action has eight elements, acting by addition of these masks.
Four plaquette supports and two winding supports are
\[
\begin{array}{c|l}
 p_0&0,1,2,5\\ p_1&0,2,3,7\\p_2&1,4,5,6\\p_3&3,4,6,7\\
 w_0&0,4\\w_1&1,3
\end{array}
\]
Parity on each support is invariant under every vertex flip, and
$p_3=p_0+p_1+p_2$ in $\mathbf F_2$. Label the configuration by
the five bits $(p_0,p_1,p_2,w_0,w_1)$, in little-endian order, hence
by an integer $i\in\{0,\ldots,31\}$. These five linear forms are
independent: in a zero linear combination, link coordinates $6,7,2,4,3$
successively eliminate the coefficients of $p_2,p_1,p_0,w_0,w_1$.
The label map has rank five and kernel of dimension three. Its kernel
therefore equals the vertex-flip space. Labels are precisely the gauge
orbits, rather than just a selection of invariant observables.

Let $\beta=(\log3)/2$. Independent temporal link increments have
same/flip probabilities $3/4,1/4$, since their Wilson odds are
$e^{2\beta}=3$. A microscopic increment $u$ has probability
$3^{8-|u|}/4^8$. Its label distribution has the integer numerators
\[
\begin{array}{rrrrrrrr}
7048&1800&3000&3000&3000&3000&1800&2952\\
1800&1800&1080&3000&3000&1080&1800&1800\\
1800&1800&3000&1080&1080&3000&1800&1800\\
648&1800&1080&1080&1080&1080&1800&648
\end{array}
\]
in label order $0$ through $31$. Denote these numbers by $n_i$.
They are obtained by summing $3^{8-|u|}$ over the eight configurations
in each label fiber; their sum is $65536=4^8$. The actual kinetic
matrix on the quotient is consequently
\[
 K_{ij}=n_{i\mathbin\oplus j}/65536.
\]
Linearity of the label map proves this convolution formula directly
from the microscopic independent increments. It is not an assumed
quotient operator. Every entry is positive, and $K$ is symmetric and
stochastic, with uniform stationary law.

The sum of the four microscopic plaquette signs is
\[
 S(i)=\sum_{r=0}^3(-1)^{p_r(i)}.
\]
Put $o(i)=(0,2,1,1,1,1,2,3)_{i\bmod8}$, so
$S(i)=(4,0,0,-4)_{o(i)}$. The spatial half weight at the very same
coupling is
\[
 H_i=e^{\beta S(i)/2}=(3,1,1,1/3)_{o(i)},\qquad
 T_{ij}=H_iK_{ij}H_j.
\]
All eight temporal links and all four spatial plaquettes thus have one
common $\beta$. There is no orbit-dependent coupling.

\subsection*{A positive ground function with exact algebraic data}
For every function $v$ of the four orbit types, direct aggregation of
the displayed $32$-state matrix gives
\[
 \sum_jT_{ij}v_{o(j)}=(Bv)_{o(i)},\qquad
 B=\begin{pmatrix}
3177/2048&765/512&675/1024&225/2048\\
765/2048&257/512&255/1024&85/2048\\
675/2048&255/512&289/1024&75/2048\\
225/2048&85/512&75/1024&353/18432
\end{pmatrix}.
\]
This identity for all $i$ and all four coordinate basis vectors is one
of the exact finite checks. Define
\begin{align*}
 P(x)&=589824x^4-1388800x^3+284192x^2-5425x+9,\\
 v(x)&=\begin{pmatrix}
4718592x^3-3790592x^2+108928x-353\\
1762560x^2-122400x+765\\
1555200x^2+86400x-675\\
518400x^2-2025
\end{pmatrix}.
\end{align*}
Exact polynomial arithmetic gives
\[
 (B-xI)v(x)=(-8P(x),0,0,0)^T.
\]
Also $P(2)<0<P(9/4)$, so the intermediate value theorem supplies
$\lambda\in[2,9/4]$ with $P(\lambda)=0$. All coordinates of $v(x)$
are strictly positive for $x\geq2$. For the first, write
$x^2(4718592x-3790592)+108928x-353$; for the second use
$x(1762560x-122400)+765$. Both are positive on that interval, as are
the last two by their displayed coefficients and $x^2\geq4$.
Thus $g_i=v_{o(i)}(\lambda)>0$ is a genuine constructed eigenfunction
of $T$, with eigenvalue $\lambda>0$.

It follows that
\[
 Q_{ij}=\frac{T_{ij}g_j}{\lambda g_i},\qquad
 \pi_i=\frac{g_i^2}{\sum_jg_j^2}
\]
define a positive stochastic matrix and a probability law. The row sum
identity is exactly $Tg=\lambda g$. Symmetry of $T$ gives
$\pi_iQ_{ij}=\pi_jQ_{ji}$, hence stationarity after summing in $i$.
This supplies the normalized stationary model without postulating a
vacuum or relying on a numerical eigenvector.

\subsection*{The sharp bare gap and its strict deterioration}
For $a\in\mathbf F_2^5$, the Walsh character
$\chi_a(i)=(-1)^{a\cdot i}$ is an eigenfunction of $K$. Pulling it
back to the links gives parity on the symmetric difference $R_a$ of
the five support rows. Since a single noise bit has sign mean $1/2$,
independence gives the exact eigenvalue $2^{-|R_a|}$. The rows are
independent, so $R_a$ is nonempty when $a\ne0$. Every row has even
incidence at each vertex; this is preserved by symmetric difference.
A single link has two odd endpoints, and therefore $|R_a|\geq2$.
The winding row $\{0,4\}$ attains two. The 32 characters are a full
orthonormal basis for uniform $L^2$: their pairwise averages factor
over the five independent uniform bits, giving $1$ on equal indices
and $0$ otherwise. Expansion in this basis proves, for every
uniformly centered $f$,
\[
 \langle f,Kf\rangle_{\rm unif}\leq\tfrac14\|f\|_{\rm unif}^2.
\]
Equality holds at a nonzero winding character. Thus the exact bare
Poincar\'e gap is $3/4$, an all-function statement rather than a
test of a few modes. This support/parity proof is a human
reconstruction of the library's complete integer character certificate.

Now use the actual gauge-invariant probe
\[
 f_i=\begin{cases}1,&i=0,16,\\-1,&i=8,24,\\0,&\text{otherwise}.
 \end{cases}
\]
It is the first winding sign restricted to configurations with all
plaquettes positive. The four supported states have identical $g_i$
and $H_i=3$, so $\E_\pi f=0$ and $\E_\pi f^2>0$. On this support,
\[
 \sum_jK_{ij}f_j=\frac{25}{256}f_i,
 \qquad (Qf)_i=\frac{225}{256\lambda}f_i.
\]
For instance the first equality at $i=0$ uses the exact numerator
$7048-1800+1800-648=6400$. The other three equalities follow by
the same exclusive-or translates. Only the supported states enter
the lag pairing, so
\begin{equation}\label{eq:z2rayleigh}
 \frac{\langle f,Qf\rangle_\pi}{\|f\|_\pi^2}
 =\frac{225}{256\lambda}\geq\frac{25}{64}>\frac14.
\end{equation}
Consequently the true stationary Poincar\'e gap is at most $39/64$,
strictly less than the sharp bare gap $3/4$. Every comparison asserting
the latter lower bound unchanged for all $\pi$-centered functions is
false in this actual homogeneous finite model. More quantitatively,
any bound $(3/4-\varepsilon)\|h\|_\pi^2\leq
\langle h,(I-Q)h\rangle_\pi$ for all centered $h$ must pay
$\varepsilon\geq9/64$, by substituting this probe into
\eqref{eq:z2rayleigh}. The probe itself has bare Rayleigh quotient
$25/256$, not the maximal bare quotient $1/4$; these are different
comparisons.

Strict positivity still gives fixed-model mixing: with $\eta=\min_{i,j}Q_{ij}/\pi_j>0$, row sums give $\eta\leq1$ and $Q\geq\eta\,1\otimes\pi$. If $\eta<1$, the remaining normalized Markov matrix preserves $\pi$ and contracts $L^2(\pi)$ by Jensen, so the centered operator norm is at most $1-\eta$; if $\eta=1$ it is zero.
The example disproves preservation of a particular kinetic gap, not
existence of every gap. It neither settles the equally weighted
seven-link $\SU$ regime nor supplies a volume-uniform or physical
four-dimensional conclusion. Positive ground normalization and Walsh
calculus are classical. The exact finite comparison has unresolved
priority, and no novelty is inferred from its finite verification.
\par\smallskip\begingroup\small\raggedright
\textit{Source correspondence:} \code{FiniteWilsonTorusNoise},
\code{FiniteWilsonTorusGround}, \code{FiniteWilsonTorusStationary},
\code{FiniteWilsonTorusCharacters}, \code{FiniteWilsonTorusGap}.
\par\endgroup


\subsection*{An elementary plaquette observable also slows}
The preceding probe contains a winding sign. Slowing also occurs for
a microscopic observable specified entirely by elementary plaquettes:
the score $S(i)=\sum_{r=0}^3(-1)^{p_r(i)}$ already used in the
spatial weight. Its centering must be performed separately in each
stationary law. Write
\[
 \widehat S=S-\mathbb E_\pi S,
 \quad V=\mathbb E_\pi\widehat S^2,
 \quad L=\langle\widehat S,Q\widehat S\rangle_\pi.
\]
We prove
\begin{equation}\label{eq:plaquetteslowing}
 \frac{\langle S,KS\rangle_{\rm unif}}{\|S\|_{\rm unif}^2}
 =\frac1{16}<\frac{L}{V},\qquad V>0.
\end{equation}
This compares normalized one-step centered covariances for the same
microscopic score in two genuine stationary laws. It does not claim
that the score is an eigenfunction of the weighted operator, nor
identify its quotient with the full stationary spectral gap.

For the bare law, each elementary plaquette is parity on four distinct
links. A noise increment has sign mean $1/2$, so independence gives
$K(-1)^{p_r}=2^{-4}(-1)^{p_r}$. The four corresponding nonconstant
Walsh characters are distinct and orthogonal in the uniform quotient
law. Therefore $\mathbb E_{\rm unif}S=0$,
$\mathbb E_{\rm unif}S^2=4$ and
$\langle S,KS\rangle_{\rm unif}=1/4$, proving the first equality.

Put $v_j=v_j(\lambda)$ for the positive polynomial vector above and
set
\begin{align*}
 Z&=v_0^2+4v_1^2+2v_2^2+v_3^2,\\
 C&=\frac{3177}{2048}v_0^2
    -\frac{450}{2048}v_0v_3+\frac{353}{18432}v_3^2.
\end{align*}
The four orbit multiplicities in the 32-state space are $(4,16,8,4)$,
so the full stationary normalizer is $4Z>0$. Since the score is
$(4,0,0,-4)$ on these orbits, direct summation gives
\begin{equation}\label{eq:plaquettemoments}
 m:=\mathbb E_\pi S=\frac{4(v_0^2-v_3^2)}Z,\qquad
 M:=\mathbb E_\pi S^2=\frac{16(v_0^2+v_3^2)}Z.
\end{equation}
For the uncentered lag, use the genuine ground transform:
\[
 \ell:=\langle S,QS\rangle_\pi
 =\frac{\sum_{i,j}g_iS(i)T_{ij}g_jS(j)}{4\lambda Z}
 =\frac{16C}{\lambda Z}.
\]
Only orbit types zero and three contribute. Their four-orbit entries
$B_{00},B_{03}=B_{30},B_{33}$ produce exactly the displayed $C$.
Stochasticity and stationarity then imply
$V=M-m^2$ and $L=\ell-m^2$: expand the centered products and use
$Q1=1$ and $\mathbb E_\pi QS=\mathbb E_\pi S$.
Every $\pi_i$ is positive and the score takes both $4$ and $-4$,
so $V$ is strictly positive.

Here is a rational certificate for the strict inequality, avoiding
numerical ground-state approximations. Define the polynomials $Z(x)$
and $C(x)$ by replacing each $v_j$ in their formulas by $v_j(x)$, and
put
\begin{align*}
 D(x)={}&256Z(x)C(x)
 -16xZ(x)\bigl(v_0(x)^2+v_3(x)^2\bigr)\\
 &-240x\bigl(v_0(x)^2-v_3(x)^2\bigr)^2.
\end{align*}
Modulo the displayed polynomial $P(x)$, exact rational reduction gives
\begin{equation}\label{eq:plaquetteremainder}
 D(x)\equiv\frac{73153125}{65536}\,R(x-2)\pmod{P(x)},
\end{equation}
where $R(y)=a_0+a_1y+a_2y^2+a_3y^3$ has coefficients
\[
\begin{array}{c|r}
0&3221027340138603340546114739\\
1&5030651367929054084455750651\\
2&2615855257226141153611982240\\
3&452893609599424244451318528
\end{array}
\]
For reproducibility, the reduction uses only repeated substitution of
\[
 x^4=\frac{1388800x^3-284192x^2+5425x-9}{589824}
 \quad\pmod{P(x)}
\]
after expanding the explicit polynomials $v_j,Z,C,D$, followed by the
binomial substitution $x=2+y$. Thus \eqref{eq:plaquetteremainder}
is a finite rational polynomial identity, not an approximate sign test.
Its coefficient arrays are independently checked with exact division
in the accompanying script. Every $a_j$ is positive, and
$\lambda\geq2$ and $P(\lambda)=0$, so $D(\lambda)>0$.

Finally, the exact moment formulas give
\[
 (16\ell-M-15m^2)\lambda Z^2=D(\lambda)>0.
\]
Since $\lambda Z^2>0$, this says
$\ell-m^2>(M-m^2)/16$, and division by $V>0$ proves
\eqref{eq:plaquetteslowing}. This proof exposes the actual centering
and algebraic remainder in the formal module and declarations
listed below. The positive shifted remainder
is a human reorganization of the same formal rational certificate.
It is a fixed finite-model calculation with unresolved priority;
no volume-uniform local correlation bound or continuum conclusion is
inferred from it.
\par\smallskip\begingroup\small\raggedright
\textit{Source correspondence:} \code{FiniteWilsonPlaquetteSlowing};
\code{centered_score_slowing},
\code{centered_score_variance_pos},
\code{plaquette_covariance_increases}.
\par\endgroup



\section{Physical scaling and the exact positive spectral bridge}
\label{sec:scalingspectral}
The finite comparisons above have no lattice-spacing parameter.
The following separate conditional results explain what additional
information a scaling argument must retain. They assert nothing about
the existence of a particular continuum Yang--Mills measure.

\subsection*{Uniform block decay can force a zero physical limit}
Suppose a real connected-correlation sequence satisfies
\begin{equation}\label{eq:blockdecayscaling}
 |c_n|\leq Cq^n,\qquad C\geq0,\quad0\leq q<1.
\end{equation}
Then $c_n\to0$, since $q^n\to0$ and its positive and negative
envelopes both tend to zero. Consequently this sequence cannot also
converge to a nonzero real limit. The same argument works if the
exponent is an integer block count $N_j\to\infty$: composition of
the geometric limit gives $Cq^{N_j}\to0$. This reindexing is an
elementary human consequence of the mapped sequential theorem.

For example, if a positive fixed physical separation is represented
by a growing number of lattice blocks as the spacing tends to zero,
a spacing-independent $q<1$ and amplitude $C$ force the corresponding
unrenormalized correlations to vanish. This conclusion depends on
both uniformities. It is not a conclusion that every possible field
normalization has zero correlations, or that the finite theorem
already provides either of the uniform hypotheses.

There is also an exact obstruction to one proposed normalization.
Suppose $c_n>0$ and the variance of one raw centered observable is
$v_n\geq v>0$. Multiplying both observables by $c_n^{-1/2}$ makes
their cross covariance one, and makes that observable's variance
$v_n/c_n$. These variances cannot have a finite common upper bound:
if $v_n/c_n\leq M$, positivity gives
$0<v\leq v_n\leq Mc_n\to0$, a contradiction.
The lower bound on the raw variance is essential to this argument;
it does not rule out normalizations whose raw variances also vanish.
\par\noindent\code{YMContinuumCorrelationScaling}:
\par\noindent\code{uniform_block_bound_forces_zero},
\par\noindent\code{no_nonzero_limit_with_uniform_block_bound},
\par\noindent\code{no_uniform_variance_after_unit_covariance}.

For comparison of units, if a single step has physical length
$a>0$ and contraction $0<q_a<1$, its exponential decay rate in
physical units is $-\log(q_a)/a$. A fixed $q_a=q<1$ gives an
unbounded rate as $a\downarrow0$, consistently with the zero-limit
conclusion above under the uniform amplitude bound. A finite rate
$M>0$ instead has the illustrative scaling $q_a=e^{-Ma}\to1$;
then $q_a^{N_a}=e^{-MaN_a}\to e^{-Mr}$ whenever $aN_a\to r$.
These are human arithmetic illustrations of units, not constructed
lattice families or renormalization theorems for our finite model.


\subsection*{An exact weak-coupling error budget and its block scale}
A separate reference expression from the audited manuscripts is
\[
 r_\beta=\exp\left(-\frac{\ell}{A\beta}\right),
             \qquad A,\ell,\beta>0.
\]
It is strictly below one for each $\beta$, but tends to one as
$\beta\to\infty$. Thus no fixed $q<1$ bounds it for all positive
$\beta$. Adding nonnegative error terms does not restore such a
uniform bound. This is a statement about the displayed reference
factor, which has not been identified with an optimal contraction
of a continuum physical transfer operator.

The amount of available room below one is at most inverse coupling.
The elementary inequality $e^{-x}\geq1-x$ gives
\[
 1-r_\beta\leq\frac{\ell}{A\beta},\qquad
 r_\beta+e_1+e_2<1
       \ \Longrightarrow\ \beta(e_1+e_2)<\frac{\ell}{A}.
\tag{\ref{eq:blockdecayscaling}$'$}\label{eq:referencebudget}
\]
For example, the proposed certified factor
$r_\beta+c/\sqrt\beta$, with fixed $c>0$, is at least one for
all sufficiently large $\beta$. Otherwise the necessary inequality
above would give $c\sqrt\beta<\ell/A$, contradicting divergence
of its left side. The same conclusion holds for an actual error
sum proved to be eventually at least $c/\sqrt\beta$.
An upper bound of order $\beta^{-1/2}$ alone does not establish
that lower bound on the true error; failure of that proposed upper
certificate is not failure of all possible contraction estimates.

Nor does $r_\beta\to1$ rule out decay over a changing number of
steps. For integers $n\geq1$, $k\geq0$, exact exponential arithmetic
gives
\[
 r_n^n=e^{-\ell/A},\qquad
 r_n^{nk}=e^{-(\ell/A)k}.
\]
Consequently for every sequence of positive integers $n(k)$,
$r_{n(k)}^{n(k)k}\to0$ as $k\to\infty$. These exact identities
separate a degenerating single-step estimate from its block scale.
They concern the reference expression itself; no accumulation bound
for other errors or identification with a physical length is supplied.
\par\noindent\code{YMBlockRateDegeneration}:
\par\noindent\code{referenceFactor_tendsto_one}, \code{referenceFactor_macro_steps},
\par\noindent\code{referenceFactor_varying_macro_decay}.
\par\noindent\code{YMBlockErrorBudget}:
\par\noindent\code{strict_factor_requires_inverse_coupling_error},
\par\noindent\code{fixed_sqrt_error_eventually_blocks_strictness}.

\subsection*{Exponential Laplace decay excludes low positive spectral mass}
Let $\mu$ be a finite positive Borel measure on $\R$. It is a
positive measure, even though its variable is allowed to range over
all real energies. Assume for every integer $n\geq0$ that
\[
 L(n)=\int_\R e^{-n\lambda}\,d\mu(\lambda)<\infty,
 \qquad L(n)\leq Ce^{-mn},
\]
where $m,C\in\R$. These are actual integrability and upper-bound
hypotheses. They imply
\begin{equation}\label{eq:laplacespectralgap}
 \mu(( -\infty,m))=0.
\end{equation}
Indeed, fix $b<m$. For $\lambda<b$ and $n\geq0$ we have
$e^{-n\lambda}\geq e^{-nb}$, hence positivity yields
\[
 e^{-nb}\mu(( -\infty,b))\leq L(n),\qquad
 0\leq\mu(( -\infty,b))\leq C e^{-(m-b)n}.
\]
Taking $n\to\infty$ gives zero mass below $b$.
Finally
$(-\infty,m)=\bigcup_{k\geq1}(-\infty,m-1/k)$,
so countable subadditivity gives \eqref{eq:laplacespectralgap}.
Mass at the endpoint $m$ is not excluded. Only integer Laplace
times were required. No cancellation argument for signed spectral
weights is being used.

Conversely, if \eqref{eq:laplacespectralgap} holds for this positive
finite measure, then for every real $t\geq0$ the integrand is
$\mu$-almost everywhere at most $e^{-mt}$, so it is integrable and
\[
 0\leq\int e^{-t\lambda}\,d\mu(\lambda)
       \leq\mu(\R)e^{-mt}.
\]
This last converse is a labelled human elementary consequence;
the named formal bridge proves the exclusion direction with its
explicit integral hypotheses.
\par\noindent\code{YMSpectralLaplaceBridge}:
\par\noindent\code{interval_lower}, \code{no_mass_below}, \code{no_mass_below_gap}.

To apply the bridge to a physical Hamiltonian, one must first have
its positive spectral measure and identify the correlation with the
displayed Laplace integral. A decay bound for one probe then controls
that probe's measure. It does not automatically give an operator-wide
lower spectral bound. An explicit human finite-dimensional example is
\[
 H=\operatorname{diag}(0,1,2)\text{ on }\R^3,\qquad
 \Omega=e_0,\quad f=e_2.
\]
The centered probe $f\perp\Omega$ has
$\langle f,e^{-tH}f\rangle=e^{-2t}$ and spectral measure $\delta_2$,
yet $e_1\perp\Omega$ has energy one. Thus this probe alone cannot
establish a lower spectral bound of two on the centered sector.
For a specified self-adjoint Hamiltonian with a spectral resolution,
one possible closing hypothesis is the same exclusion below $m$
for a dense set of centered vectors. Each corresponding spectral
projection then annihilates that dense set, and its boundedness
makes it zero on the whole centered sector. This is a conditional
human spectral-theorem consequence, not a construction of such a
Hamiltonian, spectral resolution or dense physical probe family.

Jaffe--Witten's official problem formulation requires a nontrivial
four-dimensional quantum theory as well as a positive gap; its
preceding definition distinguishes the vacuum from positive-energy
spectrum and requires a finite mass~\cite{JW}. Our elementary bounds
above isolate two necessary pieces of a proposed argument. They do
not supply the missing theory, observable normalization, spectral
identification or volume/spacing-uniform estimates. Nor does a
failure of those proposed pieces refute the physical mass gap.


\section{Two elementary clustering inferences that fail}
The following two checks concern different models and are independent
of the seven-link trial comparison. They identify missing hypotheses
in stronger clustering arguments without supplying a continuum theory.

\subsection*{Reflected positivity does not preserve translated norms}
Consider the scalar kernel $C(u)=e^{-|u|}$ and its reflected
positive-time kernel
\[
 K(s,t)=C(s+t)=e^{-s}e^{-t}\qquad(s,t\geq0).
\]
For every finite list of nonnegative times and real coefficients,
\[
 \sum_{i,j}c_i c_jK(t_i,t_j)
     =\left(\sum_i c_i e^{-t_i}\right)^2\geq0.
\]
This proves reflection positivity at this scalar two-point level
directly. The semidefinite norm of a formal positive-time test
$v=\sum_i c_i e_{t_i}$ is
$\|v\|_K=|\sum_i c_i e^{-t_i}|$. Quotienting its null space
identifies the resulting real Hilbert space with $\R$.
Translation by $a\geq0$, $T_a e_t=e_{t+a}$, is well defined on
that quotient because
\[
 \|T_av\|_K=e^{-a}\|v\|_K.
\]
It is strictly contractive for every nonzero vector when $a>0$.
In particular $K(a,a)=e^{-2a}<K(0,0)=1$. The Hilbert quotient
and exact scalar action are human reformulations of the verified
finite Gram identity, not a full field reconstruction.

The distinction also appears in a mixed pairing. For $s>0,t\geq0$,
\[
 K(s,s+t)=e^{-(2s+t)}<e^{-t}=C(t).
\]
Simultaneously translating the arguments of the unreflected
difference kernel $C(v-u)$ leaves it unchanged. Reflecting the
first positive-time argument instead gives the sum $s+t$; the
same translation now changes that sum by twice the shift.
Thus stationarity of the difference kernel and positivity of the
reflected Gram form do not justify an equality between these two
pairings or a translated reflected-norm isometry.

This agrees with the classical transfer interpretation: L\"uscher's
primary introduction~\cite{Luscher}, equations (4)--(5), constructs
the reflected inner product and identifies positive Euclidean time
translation with $e^{-aH}$. We use that passage for context only,
not as a full reconstruction theorem for this scalar example.
No field axioms, vacuum, or Yang--Mills spectral conclusion are
deduced from the rank-one kernel.

\subsection*{A nonzero actual gauge covariance needs observable normalization}
For a separate two-loop probability model, take product Haar law
$\mu=dg\,dh$ on $\SU^2$, and set
\[
 f(A,B)=a(A),\quad g(A,B)=a(B),\quad x=fg,\qquad
 d\mu_r=\frac{e^{rx}}{\int e^{rx}\,d\mu}\,d\mu.
\]
These functions and the tilted law are invariant under separate
conjugation of the two loops, and $|f|,|g|,|x|\leq1$.
Let $c(r)=\operatorname{Cov}_{\mu_r}(f,g)$.
The Haar moments from Section 2 and independence give
\[
 \E_\mu f=\E_\mu g=\E_\mu x=0,\qquad
 \E_\mu x^2=(\E a^2)^2=\frac1{16}.
\]
For bounded $u$, differentiation of the normalized exponential
integrals gives
\[
 \left.\frac d{dr}\E_{\mu_r}u\right|_{r=0}
   =\E_\mu(ux)-\E_\mu u\,\E_\mu x.
\]
This differentiation is justified on $|r|\leq1$ by the constant
majorant $e\|u\|_\infty$, and the positive denominator has value
one at zero. Applying the product rule to the covariance therefore
gives $c(0)=0$ and $c'(0)=1/16$.
Consequently every $\varepsilon>0$ contains some
$0<r<\varepsilon$ with $c(r)>0$.

Fix any such $r$ and any real $m,t$. There is no finite real $C$
such that
\begin{equation}\label{eq:scaled-cluster-impossible}
 |\operatorname{Cov}_{\mu_r}(b f,g)|\leq C e^{-mt}
                       \qquad\hbox{for every real }b.
\end{equation}
Indeed covariance is linear in $b$, so the left side is
$|b|c(r)$. The case $b=0$ would force $D=Ce^{-mt}\geq0$;
then $b=(D+1)/c(r)$ violates the bound by exactly one.
This is an actual gauge-invariant finite-law witness, not an
assumed positive covariance. It excludes only an observable-independent
prefactor on an unnormalized class: bounds retaining observable
norms or a fixed normalized class are not ruled out.
It gives no statement about separation decay or a physical mass gap.
Both mechanisms here are elementary classical obstructions, with
their precise original models and quantifiers retained.

\par\noindent\begingroup\small\raggedright
\textit{Source correspondence:} \code{YMReflectionPositiveShift},
\code{YMObservableScaleObstruction} (file \code{YMOSClusterObstruction.lean}).
\par\noindent
\code{SU2CoupledLoopResponse.corr_hasDerivAt_zero},
\code{positive_corr_near_zero}, \code{coupledLaw_gauge_invariant}.
\par\endgroup

\section{Provenance, classical context and remaining boundaries}
Every formal source used here is frozen at the verified commit
\begin{center}\code{bab688bb6140cd6316c80680095a914ea0d346db}.\end{center}
The accompanying theorem map records exact declarations, unabridged
headers, file hashes and baseline links. The original full-build and
transitive-axiom receipt, with its evidence hashes, accompanies the
draft. The publication editor has not changed a mathematical source
or run a redundant research build.

The principal sources are the shared-plaquette Haar projection and
channel modules; seven-link transfer, variational, odd-trial and spatial
separation; positive kernel powers; realized vacuum and original ground
transform; full stationary $L^2$ transfer; iteration, mixing and gauge
averaging; and the homogeneous-plaquette obstruction. The separate $\mathbf Z_2$ section includes all microscopic gauge/noise fibers, the actual 32-state transfer, four-orbit algebraic ground data, the full character certificate and the stationary gap comparison. Its checker \texttt{verify\_homogeneous\_z2.py} verifies 16 exact groups with no floating eigenvector. Its support/parity proof of the sharp bare gap is a labelled human reconstruction of those complete formal character identities.
In particular, the earlier conditional ground-data contract is
discharged by the proved positive-power construction. It is not used
as an existence axiom. Quaternion coordinate calculations, the factored
optimized margin, explicit smallness window and even/odd iteration
reorganize the verified arguments into human proofs; they are not
represented as new Lean declarations. The same exact one-weight model,
parameter dependence, centering and boundedness hypotheses are retained.

Positive-kernel existence, minorization, ground transforms and character
methods are classical. Birkhoff's primary introduction and
uniform-positivity framework \cite{Birkhoff} provide relevant antecedents;
full theorem-by-theorem historical comparison remains incomplete.
L\"uscher's primary construction \cite{Luscher} includes the physical
gauge-invariant Hilbert space and transfer matrix of lattice gauge
theory. Here the finite Markov covariance is not identified with a
physical Schwinger function or energy spectrum. Priority for the exact
optimized seven-link comparison remains unresolved after a bounded
primary-literature review.

There is no volume or lattice-spacing parameter in the theorem.
No uniform homogeneous $\SU$ coupling regime, reflection-positive continuum
construction, spectral-to-physical distance identification or
four-dimensional mass gap is supplied. Failure of a parameter-uniform
number of these finite Markov steps does not refute a physical mass gap.
The draft exposes finite-transfer mechanisms and exact hypotheses for
their reuse. The active research campaign and its unresolved continuum
objectives remain separate.

\begin{thebibliography}{9}
\bibitem{Birkhoff} G. Birkhoff,
\emph{Extensions of Jentzsch's theorem}, Transactions of the American
Mathematical Society \textbf{85} (1957), 219--227.
\href{https://doi.org/10.1090/S0002-9947-1957-0087058-6}
{doi:10.1090/S0002-9947-1957-0087058-6}.
\bibitem{Luscher} M. L\"uscher,
\emph{Construction of a selfadjoint, strictly positive transfer matrix
for Euclidean lattice gauge theories}, Communications in Mathematical
Physics \textbf{54} (1977), 283--292.
\href{https://doi.org/10.1007/BF01614090}{doi:10.1007/BF01614090}.
\href{https://bib-pubdb1.desy.de/record/396349/files/7611148.pdf}
{Primary DESY preprint}.

\bibitem{JW} A. Jaffe and E. Witten,
\emph{Quantum Yang--Mills theory}, official Clay Mathematics Institute
problem description, Section 4, printed page 6.
\href{https://www.claymath.org/wp-content/uploads/2022/06/yangmills.pdf}{Primary problem statement}.
Definition comparison only.
\end{thebibliography}
\end{document}
